\newcommand{\TrainingPanel}[2]{%
\begin{tikzpicture}
\begin{axis}[width=\linewidth,height=5.6cm,title={#2},
 xlabel={刺激序号},ylabel={aBN1 脉冲数},xmin=0.7,xmax=12.3,ymin=-0.3,ymax=11,
 xtick={1,3,6,9,12},ytick={0,2,4,6,8,10},grid=major,grid style={black!10},
 tick label style={font=\footnotesize},label style={font=\small},title style={font=\small\bfseries}]
\addplot[name path=staticlow,draw=none,forget plot] table[x=trial,y=low]{data/train_static_#1.csv};
\addplot[name path=statichigh,draw=none,forget plot] table[x=trial,y=high]{data/train_static_#1.csv};
\addplot[Gray!15,forget plot] fill between[of=staticlow and statichigh];
\addplot[name path=stdlow,draw=none,forget plot] table[x=trial,y=low]{data/train_std_#1.csv};
\addplot[name path=stdhigh,draw=none,forget plot] table[x=trial,y=high]{data/train_std_#1.csv};
\addplot[Blue!15,forget plot] fill between[of=stdlow and stdhigh];
\addplot[Gray,thick,mark=square*,mark size=1.5pt] table[x=trial,y=mean]{data/train_static_#1.csv};
\addplot[Blue,thick,mark=*,mark size=1.6pt] table[x=trial,y=mean]{data/train_std_#1.csv};
\end{axis}
\end{tikzpicture}}
\begin{figure}[tbp]
\centering
\begin{minipage}{0.49\linewidth}\TrainingPanel{500}{A. 起始间隔 0.5 s}\end{minipage}\hfill
\begin{minipage}{0.49\linewidth}\TrainingPanel{2000}{B. 起始间隔 2 s}\end{minipage}
\par\smallskip{\small\textcolor{Gray}{$\blacksquare$ M0：固定连接}\qquad\textcolor{Blue}{$\bullet$ M1：局部 STD}}
\caption{重复刺激下 aBN1 的响应。每点为五个冻结输入实现的均值；阴影为最小值至最大值，不是置信区间。计数窗口为每次刺激起始后的 300 ms。两幅图纵轴相同。M0 逐次恒定；M1 在短间隔下强烈递减，在长间隔下保留较大响应。横轴是刺激序号，不是相同长度的物理时间。}
\label{fig:training}
\end{figure}
